% ======================================================================
%  Preface
% ======================================================================

\chapter*{Preface}
\addcontentsline{toc}{chapter}{Preface}
\markboth{Preface}{Preface}

This book is about light that refuses to turn back.

In an ordinary optical waveguide, every imperfection---a bend, a
scratch, a change in refractive index---scatters some light backward.
Over long distances, backward scattering accumulates and degrades any
signal carried by the light.  This is not an engineering failure; it is
a consequence of time-reversal symmetry, which guarantees that every
forward-propagating mode has a backward partner at the same frequency.

In 2005, Haldane and Raghu proposed a way to break this symmetry in
photonic crystals, creating channels where light can propagate in only
one direction.  These ``one-way waveguides'' are the photonic analogs of
the chiral edge states in the quantum Hall effect.  In 2009, Wang,
Chong, Joannopoulos, and Solja\v{c}i\'{c} demonstrated the effect
experimentally: microwave radiation travelling along the edge of a
gyromagnetic photonic crystal with a forward-to-backward ratio
exceeding $10^5$.

The physics behind this phenomenon is \emph{topological}.  Each
electromagnetic band in a periodic medium carries a global invariant
---the Chern number---that counts how many times the Berry phase of the
Bloch modes winds as the crystal momentum sweeps the Brillouin zone.
When two media with different Chern numbers are joined, unidirectional
edge modes necessarily appear at the interface.  The number and
direction of these modes are fixed by topology and cannot be removed by
any perturbation that preserves the bulk band gap.

\paragraph{Why this book.}

Topological photonics has grown rapidly since 2005.  Excellent reviews
exist, notably those by Lu, Joannopoulos, and Solja\v{c}i\'{c} (2014),
Khanikaev and Shvets (2017), and Ozawa et~al.\ (2019).  Yet there is no
monograph that develops the subject \emph{from Maxwell's equations},
with the care and completeness expected of a graduate-level text in
mathematical physics.

Most treatments borrow the formalism of electronic topological
insulators---tight-binding Hamiltonians, Bloch electrons, Kramers
degeneracy---and graft photonic examples onto it.  This approach is
efficient for researchers already fluent in condensed-matter topology,
but it obscures what is specific to electromagnetism: the six-component
field structure, the material-dependent energy inner product, the
bosonic time reversal with $\Top^2=+1$, and the role of dispersive and
bianisotropic media.

This book takes a \emph{Maxwell-first} approach.  Every concept is
derived from the macroscopic Maxwell equations and the electromagnetic
energy inner product, not imported from condensed-matter theory by
analogy.  Electronic parallels are mentioned for orientation, but they
never replace the electromagnetic derivations.

\paragraph{Structure.}

The book is organised in five parts.

\textbf{Part~I} (Chapters 1--5) builds the foundation: the Maxwell
equations as a first-order dynamical system, the electromagnetic energy
inner product, normal modes, Bloch theory for periodic media, and the
Berry connection, curvature, and Chern number for Maxwell bands.

\textbf{Part~II} (Chapters 6--10) develops the theory of photonic Chern
insulators: Dirac cones from lattice symmetry, time-reversal breaking
by gyrotropic media, the one-way interface mode at a domain wall,
the Wang microwave experiment, and numerical methods for computing
Chern numbers.

\textbf{Part~III} (Chapters 11--14) covers other two-dimensional
topological phases: reciprocal designs with pseudospin, valley-Hall
photonic crystals, Floquet (time-modulated) systems, and
one-dimensional topology including the Zak phase and the SSH model.

\textbf{Part~IV} (Chapters 15--16) extends the framework to three
dimensions (Weyl points and surface arcs) and synthetic dimensions
(internal degrees of freedom as extra lattice coordinates).

\textbf{Part~V} (Chapters 17--18) addresses non-Hermitian topology
(loss, gain, exceptional points, topological lasers) and nonlinear
and quantum frontiers (topological solitons, many-body photonic
states, polaritons).

Four appendices provide reference material: vector identities and
functional analysis (A), differential forms and Chern classes (B),
numerical recipes with pseudocode (C), and material tensors with sign
conventions (D).

\paragraph{Prerequisites and audience.}

The intended reader has completed graduate courses in electrodynamics
(at the level of Jackson or Griffiths) and has some exposure to
quantum mechanics and solid-state physics.  No prior knowledge of
topology, differential geometry, or condensed-matter topological
insulators is assumed; these ideas are developed as they arise,
motivated by the electromagnetic problems that require them.

\paragraph{Exercises.}

Each chapter ends with exercises ranging from straightforward
verifications to open-ended explorations.  Working through these
exercises is essential for developing the physical and mathematical
intuition needed to apply topological ideas to new photonic systems.

\paragraph{Conventions.}

We use SI units throughout, with the time-harmonic convention
$\ee^{-\ii\omega t}$ and $\omega>0$ for physical frequencies.  The
complete notation is summarised in the next chapter and in Appendix~D.

\paragraph{Acknowledgements.}

This work builds on the contributions of many researchers who have
shaped topological photonics into a vibrant field.  Special
acknowledgement is due to the pioneering theoretical work of Haldane
and Raghu, the experimental breakthroughs of Wang, Chong, Joannopoulos,
and Solja\v{c}i\'{c}, and the comprehensive reviews that have made the
field accessible to a broad audience.

\vspace{1cm}

\hfill\emph{The Authors}

\hfill\emph{2026}
